\section{The Cooper Pair Interface}

The electron reaches the Cooper-pair boundary as a counted loop, not as an
open traveler. It carries the native negative charge, the strain face of the
second difference, and the phase face that can bounce under a tilted
comparison. The boundary now asks a new question. Can charge pass through a
region whose admissible transport cancels internal antisymmetric strain
without losing the loop count that made charge readable?

The answer is an interface. A superconducting region is not treated here as a
private interior that swallows the grammar. It is a boundary discipline. It
permits paired transport in which two charge carriers are read together so
that their internal antisymmetric residue cancels. The cancellation is real
for the internal strain of the pair. It is not permission to forget charge,
phase, or boundary cost. The grammar interfaces the pair with the loop
ledger rather than fusing the pair into an unmeasured interior.

This is why the unit is \(2e\). The superconducting boundary does not admit
solitary open charge as the primitive transport. It admits the paired
closure. Two electron charges are carried as one coherent boundary unit. The
unit is counted as \(2e\) because the boundary has paired the loop readings
before transport is accepted. The count is still a loop count; it has simply
been interfaced through a pair rather than through a single open path.

The pairing also explains the zero-strain temptation and its limit. Inside
the pair, the antisymmetric component of transport can cancel. That is the
strain-free face of coherent superconducting motion. But a strain-free
internal update is not a cost-free grammar. The phase of the pair still has
to be carried. The boundary must preserve coherent orientation across the
paired unit, and that preservation charges phase cost to the ledger.

The Meissner illustration shows the boundary form. A zero-strain region
cannot simply admit every external curvature through its interior, because
doing so would reintroduce the antisymmetric residue the paired transport
has excluded. The field is therefore read as diverted around the region.
The example is not the spine. It illustrates the grammar's rule: a boundary
that enforces strain-free paired transport also forbids incompatible
curvature from becoming an interior reading.

The gauge count gives the pair its finite audit. In the native reading, the
electron corner never cranks. Every gate pools with the matter count. At the
Cooper-pair interface, the pair remains close to that corner, but it admits
one selecting turn through the gauge. Thirty-five gates may still pool while
one gate selects the positive orientation needed by the paired boundary
reading. The pair does not become an antimatter world. It carries exactly
one positive turn as part of closing the paired ledger.

This one-turn discipline is what keeps the boundary from overclaiming. If
every gate selected, the native matter count would be erased. If no gate
selected, the pair would never reach the tilted positive orientation needed
for the boundary interface. The grammar forces the intermediate reading: a
single positive selection carried through a gauge that otherwise remains
pooled with the native electron corner. The pair is one turn away from the
electron, not an abandonment of the electron.

The topological integer-count image belongs here as a helpful shadow. Paired
transport still enters the record by whole units. The boundary does not
measure half a coherent pair any more than the loop count measures half a
closed traversal. A pair is admitted as a pair, with its boundary charge unit
and its phase relation. If the pair breaks, the grammar has a different
record: unpaired transport, strain, and resistance return.

The superconducting example makes the same point in physical language. At
low refinement noise, correlated charge carriers can move as a symmetric
unit. Their internal antisymmetric residue vanishes under the reading that
admits the pair, so transport can proceed without dissipative strain. But the
coherent unit has a phase. To move the unit through the boundary is to carry
that phase consistently. The pair's freedom from internal tearing does not
erase the cost of keeping the boundary orientation coherent.

Thus the Cooper-pair interface brings the chapter's faces together. Charge
is counted in units of paired loop transport, \(2e\). Mass is suppressed
inside the pair where antisymmetric strain cancels, but strain remains the
grammar that says what would return if the pair failed. Phase is the cost of
coherent boundary orientation. Sameness is the quotient that lets the pair
stand as one transport unit for this reading. Four faces, one paired
interface.

The chapter stops at the boundary on purpose. It does not need to derive the
full later horizon of superconductivity. It needs to carry the electron to
the place where solitary charge is no longer the admissible unit and where
phase coherence becomes the price of transport. The grammar has forced the
meeting: loop count, second-difference strain, phase bounce, and bounded
sameness all arrive at the Cooper-pair boundary.

What crosses the boundary is therefore neither a free electron story nor a
new primitive pair story. It is a charged loop ledger interfaced through
paired closure. The boundary permits \(2e\) transport, forbids incompatible
interior strain, identifies the pair as one coherent unit, carries one
positive gauge turn, and charges phase cost for keeping the closure coherent.
